% ===================================================================
%  TABEL Nr. 1 — Die skool. — Die liseum. — Die klaskamer.
%
%  Ceci n'est PAS une transcription : c'est une traduction, faite en
%  2026, en afrikaans d'aujourd'hui. D'ou trois differences avec
%  texto/io et texto/fr, et trois seulement :
%
%   * pas de \nl ni de \cc — la traduction ne suit aucune ligne
%     imprimee, et les lignes de ce fichier ne sont que du confort ;
%   * pas de \begin{VUpage} — elle n'a ni feuillet ni folio, et la
%     page de lecture ne lui en affiche pas ;
%   * le gras se pose sur le TERME seul, ponctuation dehors.
%
%  Tout le reste est identique, et doit l'etre : les cles « %%K » sont
%  celles de texto/io, macro pour macro pour l'apparat, et les renvois
%  \textsuperscript{(n)} visent les MEMES objets de la planche.
%
%  UNE COLONNE HORS DES DEUX LISTES, comme le polonais : ni les
%  vingt-neuf d'Ethnologue, ni les dix-sept pays de l'ido. Elle est la
%  parce qu'on l'a demandee.
%
%  ET SA VOISINE EST DANS LE MEME DOSSIER. texto/nl est la langue dont
%  celle-ci sort : l'afrikaans est le neerlandais des colons du Cap,
%  trois siecles plus tard. C'est la situation du bhojpouri devant le
%  hindi — a une difference pres, decisive : ICI LA DEFENSE SE JOUE
%  AUSSI A LA LETTRE.
%
%   * le digramme « ij » n'existe pas : zijn, tijd, wij, vrij donnent
%     is, tyd, ons, vry. Deux caracteres relevent toute une classe de
%     calques ;
%   * le « z » initial non plus : zien, zitten, zon donnent sien, sit,
%     son ;
%   * « sch » au debut d'un mot est du neerlandais : school donne
%     skool.
%
%  ET LA GRAMMAIRE, ELLE, NE SE DEFEND PAS A LA LETTRE. Le verbe
%  afrikaans ne se conjugue pas — ek is, jy is, hy is, ons is —, le
%  passe se fait tout entier avec « het » plus « ge- », et la negation
%  se ferme : « nie ... nie ». Une phrase neerlandaise correcte reste
%  correcte a l'oeil ; c'est la ou l'outil ne peut rien et ou il faut
%  relire.
%
%  ON N'Y RELEVE NI « het » NI « een ». « het » est l'auxiliaire du
%  passe afrikaans, la ou le neerlandais en fait un article ; « een »
%  est le NUMERAL afrikaans, quand l'article est « 'n ». Deux mots que
%  la graphie confond, et une regle dessus crierait a chaque page.
%
%  LES PRENOMS SONT AFRIKAANSISES : Ioannes est Jan, Iakobus est
%  Jakobus, Karolus est Karel. Les noms latins de la note (*) ne
%  bougent pas : c'est d'eux que la note parle.
% ===================================================================

%%K t01-apar-0 apar
\VUsaut{2.0mm}
\VUcentre{12.6pt}{120}{VERKLARENDE HANDLEIDING}
\VUsaut{3.2mm}
\VUcentre{8.4pt}{160}{BY}
\VUsaut{2.6mm}
\VUtitre{15.6pt}{0}{die Delmas-hulptabelle}
\VUsaut{3.4mm}
\VUornamento{9pt}
\VUsaut{5.0mm}
\VUcentre{13.2pt}{90}{EERSTE REEKS}
\VUsaut{3.0mm}
\VUfilet{16mm}
\VUsaut{5.6mm}

%%K t01-apar-1 apar
\VUtitre{15.0pt}{60}{TABEL Nr. 1}
\VUsaut{1.4mm}
\VUtitre{11.4pt}{0}{Die skool, die liseum, die klaskamer.}
\VUsaut{2.2mm}
\VUfilet{20mm}
\VUsaut{6.4mm}

%%K t01-tit-1 sub
\VUpk{10.2pt}{40}{Algemene beskrywing.}
\VUsaut{3.0mm}

%%K t01-01-1 p
1. --- Hierdie tabel stel 'n \VUgras{klaskamer} in 'n \VUgras{liseum} voor. In hierdie
klaskamer staan agt \VUgras{tafels} \textsuperscript{(18)} en agt \VUgras{banke}
\textsuperscript{(32)}. Op elke bank sit drie leerlinge. Die \VUgras{onderwyser}
\textsuperscript{(7)} sit op 'n \VUgras{stoel} \textsuperscript{(8)} by sy
\VUgras{kateder} \textsuperscript{(6)}.

%%K t01-02-1 p
2. --- Links van die kateder staan 'n \VUgras{kas} \textsuperscript{(15)} met
\VUgras{glasdeure}. Regs is die \VUgras{swartbord} \textsuperscript{(1)} met die
\VUgras{spons} \textsuperscript{(2)} en die \VUgras{kryt} \textsuperscript{(3)}.

%%K t01-02-2 p
Bo die kateder hang \VUgras{muurplate} \textsuperscript{(9, 11, 12)} en 'n
\VUgras{kaart} \textsuperscript{(10)}.

%%K t01-03-1 p
3. --- Nie al die leerlinge sit op hulle \VUgras{plek} \textsuperscript{(17)} nie. Jan
\textsuperscript{(4)} staan by die bord, hy hou die spons vas en vee die
\VUgras{meetkundige figure} uit.

%%K t01-03-2 p
Pieter is naby die kateder, hy versamel die \VUgras{skryfwerk} \textsuperscript{(5)} en
bring dit na die onderwyser toe.

%%K t01-04-1 p
4. --- Hierdie onderwyser is die onderwyser van die \VUgras{lewende tale}. Hy leer die
leerlinge om vreemde tale te praat, te lees en te skryf \VUgras{(Duits, Engels, Spaans,
Italiaans, Russies} of \VUgras{Frans)}.

%%K t01-05-1 p
5. --- Die klaskamer is baie ruim en baie helder. Agterin is 'n breë \VUgras{venster}
met twee \VUgras{boligte} \textsuperscript{(78)} en die \VUgras{glasdeur} om dit te
lug en te verlig.

%%K t01-05-2 p
Hierdie klaskamer is nie koud nie. In die middel staan daar, om dit te verwarm, 'n baie
goeie \VUgras{stoof} \textsuperscript{(46)} met sy \VUgras{pyp} \textsuperscript{(45)}.

%%K t01-05-3 p
Teen die muur is \VUgras{kapstokhake} \textsuperscript{(58)} vasgemaak, daaraan hang
die pette en die jasse van die skoliere.

%%K t01-tit-2 sub
\VUsaut{5.2mm}
\VUpk{10.2pt}{40}{Die speelplek}
\VUsaut{3.6mm}

%%K t01-06-1 p
6. --- Die \VUgras{horlosie} \textsuperscript{(63)} wys vyf minute oor nege.

%%K t01-06-2 p
Die \VUgras{pouse} \textsuperscript{(76)} het pas geëindig en die les begin weer. Die
speelruimte lê op die \VUgras{binnehof} \textsuperscript{(75)}. Ons sien 'n paar
leerlinge wat speel. 'n \VUgras{Ondermeester} \textsuperscript{(74)} hou toesig oor
hulle.

%%K t01-07-1 p
7. --- Op die binnehof sien ons boonop verskillende persone, naamlik die
\VUgras{inspekteur} \textsuperscript{(73)}, die skoolhoof \VUgras{(rektor)}
\textsuperscript{(71)} en die \VUgras{sensor} \textsuperscript{(72)}.

%%K t01-07-2 p
Die inspekteur inspekteer die skole, die rektor lei die liseum, die sensor hou toesig
oor die studies.

%%K t01-07-3 p
Die vierde persoon is die \VUgras{portier} \textsuperscript{(77)}. Hy dra 'n register
onder sy arm om die afwesiges in te skryf.

%%K t01-08-1 p
8. --- Agterin die binnehof staan die \VUgras{gimnastiektoestelle}
\textsuperscript{(70)} en die werkplek vir die \VUgras{handearbeid}
\textsuperscript{(69)}.

%%K t01-08-2 p
Regs van die tabel begin ons die \VUgras{lokale} vir \VUgras{musiek} en
\VUgras{tekene} sien.

%%K t01-noto-1 noto
\VUnotes{143.2mm}{%
(*) Vir die \textit{doopname} word aanbeveel om hulle ook volgens die Latynse vorm oor
te skryf. Dit is die enigste manier om ontelbare afwykings te ontkom. Hoe sou 'n mens
trouens kies tussen \textit{Giovanni, Jean,} \textit{Johann, Jan, Hans, John, Ivan,}
ens.? Sal ons 'n aparte woordeboek vir die doopname maak? Laat ons uit die Latyn hulle
nominatief put en sê: \VUgras{Ioannes, Ioanna; Iulius, Iulia; Iakobus; Andreas; Lukas.}
(Volledige Gram.).%
}

%%K t01-tit-3 sub
\VUpk{10.2pt}{40}{Uitvoerige beskrywing.}
\VUsaut{2.4mm}

%%K t01-tit-4 sub
\VUpk{10.2pt}{40}{Die goeie leerlinge.}
\VUsaut{3.0mm}

%%K t01-09-1 p
9. --- Die leerlinge van die eerste bank regs van die kateder is \VUgras{Hendrik} (*),
\VUgras{Stefanus} en \VUgras{Karel}.

%%K t01-09-2 p
Hendrik is baie oplettend en werksaam, hy luister na die onderwyser.

Op sy tafel het hy 'n \VUgras{skrif} \textsuperscript{(28)}, 'n \VUgras{boek}
\textsuperscript{(29)}, 'n \VUgras{woordeboek} \textsuperscript{(30)} en 'n
\VUgras{pennehouer} \textsuperscript{(31)}. Sy boek het baie \VUgras{bladsye}, elke
hoofstuk bevat verskeie \VUgras{paragrawe} en elke \VUgras{reël} baie \VUgras{woorde}.

%%K t01-09-4 p
Sy buurman Stefanus \textsuperscript{(24)} is ywerig. Hy skryf in sy skrif die les vir
die volgende keer aan. Hy het 'n mooi \VUgras{handskrif}. Sy boek is toe. Ons sien net
die \VUgras{omslag} \textsuperscript{(27)}. Naby hom lê sy \VUgras{passer}
\textsuperscript{(26)}.

%%K t01-09-5 p
Karel \textsuperscript{(23)} hou sy boek in sy hand, hy maak dit oop om sy les weer te
lees. Hy het, soos elke leerling, 'n \VUgras{inkpot} \textsuperscript{(25)} voor hom. Hy
is gehoorsaam.

%%K t01-10-1 p
10. --- By die tafel langsaan is \VUgras{Lodewyk}, \VUgras{Antonius} en \VUgras{Paulus}.
Lodewyk sê sy \VUgras{les} \textsuperscript{(22)} op en antwoord op die \VUgras{vrae}
van die onderwyser. Antonius \textsuperscript{(21)} is baie onoplettend, die onderwyser
straf hom selfs somtyds. Paulus \textsuperscript{(20)} is 'n goeie \VUgras{maat},
vriendelik en hulpvaardig. Hy is baie oplettend.

%%K t01-tit-5 sub
\VUsaut{4.2mm}
\VUpk{10.2pt}{40}{Die slegte leerlinge.}
\VUsaut{3.2mm}

%%K t01-11-1 p
11. --- Agter Hendrik sit \VUgras{Georg} \textsuperscript{(38)}. Hy is nie 'n slegte
kêrel nie, maar hy is \VUgras{praterig}, onoplettend en baldadig. Sy
\VUgras{ligsinnigheid} en \VUgras{ongehoorsaamheid} veroorsaak \VUgras{steurnis} tydens
die les en dikwels verdien hy 'n streng \VUgras{waarskuwing}. Sy \VUgras{kladskrif}
\textsuperscript{(35)} is vuil, hy \VUgras{bemors} dit met ink met sy \VUgras{pen}
\textsuperscript{(34)}, daarom gebruik hy altyd sy \VUgras{uitveër}
\textsuperscript{(36)}; sy \VUgras{boeksak} \textsuperscript{(33)} is geskeur, want hy is
baie slordig. Waarom bring hy sy \VUgras{potloodhouer} \textsuperscript{(37)} na die
klaskamer van die lewende tale toe?

%%K t01-11-2 p
Sy buurman Edmund \textsuperscript{(39)} hou nie van hom nie. Hy is weliswaar 'n bietjie
traag, maar hy is stil en rustig. Hy praat nie; maar in plaas van te luister, lees hy
liewer die \VUgras{verhale} uit sy \VUgras{leesboek}.

%%K t01-11-3 p
Die derde leerling van hierdie bank is \VUgras{Lodewyk} \textsuperscript{(40)}. Hy is
vlytig. Hy skryf op 'n wit \VUgras{vel papier}. Voor hom het hy sy \VUgras{vloeipapier}
\textsuperscript{(41)} neergelê.

%%K t01-12-1 p
12. --- Die leerlinge van die tweede bank, links van die onderwyser, is baie jonk. Die
eerste, die klein \VUgras{Emiel}, kan nog nie baie goed skryf nie, hy bring sy
\VUgras{lei} \textsuperscript{(42)}, sy \VUgras{griffel} en sy \VUgras{spons}
\textsuperscript{(43)} saam.

%%K t01-12-2 p
Langs hom is \VUgras{August} en \VUgras{Bernard} \textsuperscript{(44)}. Hierdie een
werk goed, maar sy \VUgras{kleredrag} is baie slordig.

%%K t01-13-1 p
13. --- Agter hulle sit drie \VUgras{luiaards}. \VUgras{Albert} leer sy lesse nie.
\VUgras{Jakobus} \textsuperscript{(47)} slaap in plaas van te luister. Sy
\VUgras{luiheid} besorg hom \VUgras{verwyte} en selfs \VUgras{strawwe}. Ten slotte is
\VUgras{Christiaan} \textsuperscript{(48)} te ligsinnig. Hy \VUgras{gedra} hom sleg en
doen glad geen moeite om hom te verbeter nie.

%%K t01-14-1 p
14. --- Sy bure, daarenteen, is welopgevoed. \VUgras{Ernst} \textsuperscript{(51)} hou
van orde en reëlmaat, hy gebruik altyd sy \VUgras{liniaal} \textsuperscript{(50)} en sy
\VUgras{potlood} \textsuperscript{(49)}. Hy is 'n \VUgras{dagskolier}.

%%K t01-14-2 p
\VUgras{Frans} \textsuperscript{(52)} is 'n \VUgras{koshuisleerling}, hy dra 'n uniform,
eet en slaap in die liseum. Hy is 'n goeie \VUgras{maat}.

%%K t01-14-3 p
Maurits maak sy \VUgras{potlood} skerp met sy \VUgras{knipmes}
\textsuperscript{(56)}, hy is baie sorgsaam.

%%K t01-14-4 p
Hy bring al sy boeke saam, sy \VUgras{grammatika} \textsuperscript{(55)}, sy
\VUgras{notaboekie} \textsuperscript{(54)}, en selfs 'n \VUgras{skryfonderlegger}
\textsuperscript{(53)}. Sy \VUgras{skooltas} \textsuperscript{(57)} staan op die vloer
agter hom.

%%K t01-14-5 p
Die twee tafels agterin die klaskamer word deur \VUgras{goeie leerlinge}
\textsuperscript{(19)} beset.

%%K t01-tit-6 sub
\VUsaut{4.6mm}
\VUpk{10.2pt}{40}{Tekene en musiek.}
\VUsaut{3.4mm}

%%K t01-15-1 p
15. --- In die \VUgras{tekenkamer} is daar sierlike \VUgras{modelle} teen die muur
opgehang en 'n \VUgras{lamp} \textsuperscript{(62)} met 'n \VUgras{skerm}, wat aan die
plafon hang. Die \VUgras{onderwyser} \textsuperscript{(60)} is baie geduldig, hy
verbeter die \VUgras{tekeninge} \textsuperscript{(61)} van die leerlinge en leer hulle
om 'n \VUgras{borsbeeld} \textsuperscript{(59)} na die natuur te teken.

%%K t01-16-1 p
16. --- Die \VUgras{musieklokaal} lyk baie interessant. Daarin leer 'n mens
\VUgras{musiek} lees, die \VUgras{toonlere} \textsuperscript{(67)} solfeer wat op die
\VUgras{notebalk} geskryf is (do, re, mi, fa, sol, la, si\,=\,C. D. E. F. G. A. B.), die
\VUgras{sleutels} ken en bekoorlike \VUgras{melodieë} sing. Die musiekstukke word op die
\VUgras{lessenaar} \textsuperscript{(65)} gesit. Een van die leerlinge \VUgras{speel
klavier} \textsuperscript{(64)} en die \VUgras{musiekonderwyser} \textsuperscript{(66)}
\VUgras{slaan die maat} \textsuperscript{(68)}.

%%K t01-17-1 p
17. --- Ons begin 'n hoek van die \VUgras{werkplek} vir \VUgras{handearbeid}
\textsuperscript{(69)} sien, waar die seuns mag leer om hout te skaaf en yster te vyl.
Dit bestaan nie in alle liseums nie.

%%K t01-tit-7 sub
\VUsaut{5.0mm}
\VUpk{10.2pt}{40}{Die lokaal vir die wetenskappe.}
\VUsaut{3.6mm}

%%K t01-18-1 p
18. --- Die onderwyser van wiskunde leer die leerlinge \VUgras{meetkunde, rekenkunde,
berekening} en die \VUgras{metrieke stelsel}. Ons sien op die bord die vernaamste
meetkundige figure: die \VUgras{sirkel} \textsuperscript{(a)}, die \VUgras{vierkant}
\textsuperscript{(b)}, die \VUgras{driehoek} \textsuperscript{(c)}, die \VUgras{lyn}
\textsuperscript{(d)}.

%%K t01-18-2 p
Ons sien ook 'n \VUgras{bewerking}, dit is nie 'n \VUgras{optelling} nie, ook nie 'n
\VUgras{aftrekking} nie, ook nie 'n \VUgras{vermenigvuldiging} nie, dit is wel 'n
\VUgras{deling} \textsuperscript{(e)}.

%%K t01-19-1 p
19. --- Op die bord van die metrieke stelsel sien ons: die \VUgras{meter}
\textsuperscript{(a)}, die \VUgras{weegskaal} \textsuperscript{(b)} en haar
\VUgras{gewigte}, die \VUgras{liter} \textsuperscript{(e)}, die \VUgras{dekaliter}
\textsuperscript{(f)}, die \VUgras{desiliter} en die \VUgras{kubieke meter}
\textsuperscript{(d)}.

%%K t01-19-2 p
Die leerlinge studeer ook \VUgras{natuurwetenskap} \textsuperscript{(11)}, dierkunde
\textsuperscript{(a)}, plantkunde \textsuperscript{(b)}, geologie \textsuperscript{(c)}
en \VUgras{geskiedenis} \textsuperscript{(12)}.

%%K t01-19-3 p
In die kas is die toestelle vir \VUgras{fisika} \textsuperscript{(15)} en vir
\VUgras{chemie} \textsuperscript{(16)}.

%%K t01-19-4 p
Op die kas is: 'n \VUgras{bol} \textsuperscript{(14)}, 'n \VUgras{keël} en 'n
\VUgras{aardbol} \textsuperscript{(13)}.

%%K t01-19-5 p
Bo die tabelle hang 'n \VUgras{kaart} wat die vyf werelddele voorstel:
\VUgras{Europa} \textsuperscript{(g)}, \VUgras{Asië} \textsuperscript{(e)},
\VUgras{Afrika} \textsuperscript{(f)}, \VUgras{Amerika} \textsuperscript{(ab)} en
\VUgras{Oseanië} \textsuperscript{(h)} en die twee groot oseane, die \VUgras{Stille
Oseaan} \textsuperscript{(c)} en die \VUgras{Atlantiese Oseaan} \textsuperscript{(d)}.
\VUsaut{6.0mm}
\VUfilet{22mm}
